\documentclass[10pt,a4paper]{amsart}
\usepackage[margin=19mm]{geometry}
\usepackage{amsmath,amssymb,mathtools}
\usepackage[scr=rsfs,cal=euler]{mathalfa}
\usepackage{xfrac}
\usepackage[unicode,hidelinks,bookmarks=false]{hyperref}
\newcommand{\ie}{i.e.\ }
\newcommand{\cf}{cf.\ }
\newcommand{\resp}{respectively\ }
\newcommand{\Subsection}{\subsection}
\providecommand{\nobreakdash}{\nobreak\hbox{-}}
\setlength{\emergencystretch}{2em}
\multlinegap0pt
\usepackage{mathtools}
\usepackage{bm}
\usepackage{dsfont}
\usepackage{amssymb}
\usepackage{xcolor}
\newtheorem{lemma}{Lemma}
\title{\textbf{Maximum entropy based testing in network models: \\ ERGMs and constrained optimization}}
\author{Subhro Ghosh, Rathindra Nath Karmakar,, Samriddha Lahiry}
\date{Source: arXiv:2602.20844v2 (CC BY 4.0)}
\begin{document}
\maketitle

\noindent\emph{Source and licence.} \textbf{Maximum entropy based testing in network models: \\ ERGMs and constrained optimization}. Version arXiv:2602.20844v2 (CC BY 4.0), distributed by arXiv under the Creative Commons Attribution 4.0 International licence, http://creativecommons.org/licenses/by/4.0/, as recorded in the arXiv licence metadata for this exact version and shown on the abstract page https://arxiv.org/abs/2602.20844v2. Full licence notice: \texttt{licenses/arxiv-2602.20844-CC-BY-4.0.txt}.\par\smallskip

\noindent\textit{Editorial scope.} Counted unit: the article's lemma lambda\_cont (source0.tex lines 1064--1069). The article's complete original proof of it is delivered with it (source0.tex lines 2772--2944, one \texttt{proof} environment of 173 lines, which the article places away from the statement). No mathematics is written, repaired or re-derived: the statement and the proof are the article's own lines, in the article's own notation, and no retained line is substituted or re-flowed.  No further source-local context is needed: every cross-reference of every retained line resolves inside the two delivered spans.  Nothing was omitted from any retained span, so the package carries no omission record.  No retained line contains a citation, so the wrapper carries no bibliography and no bibliography entry.  No retained line contains a figure environment, an image-inclusion command or drawing code, so no image is delivered and the package declares no asset dependency.  Editorial formatting only, in addition: an AMS \texttt{amsart} preamble that restates the article's own declarations and macros used by the retained lines, namely the environments \texttt{lemma} on separate counters, together with the standard packages those lines need.  The retained environments print in this document as Lemma 1 in order of appearance, whereas the article prints the same items with the numbering of its own full document.  The complete native source of the article as distributed in the arXiv e-print for this exact version is bundled unchanged as \texttt{sources/195225/source0.tex}.
\begin{lemma}{\label{lambda_cont}}
   \begin{align*}
   \mathbb{E}\left[(Z_{\hat{\mu}}-h_0)^2e^{-\hat{\lambda}(Z_{\hat{\mu}}-h_0)}\vert \hat{\lambda}\right]&\xrightarrow{P} \mathbb{E}\left[(Z_{\mu}-h_0)^2e^{-\lambda^{\circ}(Z_{\mu}-h_0)}\right]\\
       \mathbf{Var}\left[(Z_{\hat{\mu}}-h_0)e^{-\hat{\lambda}(Z_{\hat{\mu}}-h_0)}\vert \hat{\lambda}\right]&\xrightarrow{P}\mathbf{Var}\left[(Z_{\mu}-h_0)e^{-\lambda^{\circ}(Z_{\mu}-h_0)}\right]\\
   \end{align*} 
\end{lemma}

\begin{proof}[Proof of Lemma \ref{lambda_cont}]
Let $h_0>0$ be a constant and define $\mu(\lambda)=h_0 e^{\lambda}$. Let
$Z_{\mu(\lambda)} \sim \mathrm{Poisson}(\mu(\lambda))$. Define the deterministic functions
\[
f(\lambda):=\mathbb{E}\Big[(Z_{\mu(\lambda)}-h_0)^2e^{-\lambda(Z_{\mu(\lambda)}-h_0)}\Big],
\qquad
g(\lambda):=\mathbf{Var}\Big((Z_{\mu(\lambda)}-h_0)e^{-\lambda(Z_{\mu(\lambda)}-h_0)}\Big).
\] As a first step we will prove that the functions $f(\lambda)$ and $g(\lambda)$ are continuous.\\

\textbf{Continuity of $f(\lambda)$}\\


We may rewrite
\[
f(\lambda)
= e^{\lambda h_0}\,\mathbb{E}\Big[(Z-h_0)^2 e^{-\lambda Z}\Big],
\qquad Z\sim \mathrm{Poisson}(\mu(\lambda)).
\]
Using the Poisson moment generating function we want to compute
\[
\mathbb{E}[(Z-h_0)^2 e^{-\lambda Z}]
= \mathbb{E}[Z^2 e^{-\lambda Z}]
-2h_0 \mathbb{E}[Z e^{-\lambda Z}]
+ h_0^2 \mathbb{E}[e^{-\lambda Z}].
\]
For a Poisson$(\mu)$ random variable we have
\[
\mathbb{E}[e^{-\lambda Z}]
= \exp\!\big(\mu(e^{-\lambda}-1)\big).
\]
Differentiating with respect to $\lambda$ yields
\[
\mathbb{E}[Z e^{-\lambda Z}]
= \mu e^{-\lambda}
\exp\!\big(\mu(e^{-\lambda}-1)\big),
\]
and a second differentiation gives
\[
\mathbb{E}[Z^2 e^{-\lambda Z}]
= \big(\mu e^{-\lambda}+\mu^2 e^{-2\lambda}\big)
\exp\!\big(\mu(e^{-\lambda}-1)\big).
\]
Substituting these expressions (with $\mu=\mu(\lambda)$) into the
previous expansion, we obtain

\[\mathbb{E}[(Z-h_0)^2 e^{-\lambda Z}]
= \exp\!\big(\mu(\lambda)(e^{-\lambda}-1)\big)
\Big[
\mu(\lambda)e^{-\lambda}
+ \mu(\lambda)^2 e^{-2\lambda}
-2h_0 \mu(\lambda)e^{-\lambda}
+ h_0^2
\Big].\]
Consequently,
\[
f(\lambda)
= e^{\lambda h_0}
\exp\!\big(\mu(\lambda)(e^{-\lambda}-1)\big)
\Big[
\mu(\lambda)e^{-\lambda}
+ \mu(\lambda)^2 e^{-2\lambda}
-2h_0 \mu(\lambda)e^{-\lambda}
+ h_0^2
\Big].
\]

Finally, since the map $\lambda\mapsto \mu(\lambda)=h_0 e^{\lambda}$ is continuous, it is easy to observe that $f(\lambda)$ is continuous in $\lambda$.\\ 

\textbf{Continuity of $g(\lambda)$}\\

Defining $Y_\lambda := (Z_{\mu(\lambda)}-h_0)\,e^{-\lambda(Z_{\mu(\lambda)}-h_0)}$, we will show that the map $
\lambda \mapsto \mathrm{Var}(Y_\lambda)
$
is continuous. It suffices to prove continuity of $\lambda\mapsto \mathbb{E}[Y_\lambda]$
and $\lambda\mapsto \mathbb{E}[Y_\lambda^2]$. We have
\[
\mathbb{E}[Y_\lambda]
=\mathbb{E}\Big[(Z-h_0)e^{-\lambda(Z-h_0)}\Big]
= e^{\lambda h_0}\,\mathbb{E}\Big[(Z-h_0)e^{-\lambda Z}\Big],
\qquad Z\sim \mathrm{Poisson}(\mu(\lambda)).
\]
Expanding gives
\[
\mathbb{E}\Big[(Z-h_0)e^{-\lambda Z}\Big]
=\mathbb{E}[Z e^{-\lambda Z}] - h_0\,\mathbb{E}[e^{-\lambda Z}].
\]
Using
\[
\mathbb{E}[e^{-\lambda Z}]=\exp\!\big(\mu(e^{-\lambda}-1)\big),
\qquad
\mathbb{E}[Z e^{-\lambda Z}]
=\mu e^{-\lambda}\exp\!\big(\mu(e^{-\lambda}-1)\big),
\]
one obtains
\[
\mathbb{E}[Y_\lambda]
= e^{\lambda h_0}\exp\!\big(\mu(\lambda)(e^{-\lambda}-1)\big)\Big(\mu(\lambda)e^{-\lambda}-h_0\Big).
\]
It is immediate that
$\lambda\mapsto \mathbb{E}[Y_\lambda]$ is continuous. Next note that
\[
Y_\lambda^2 = (Z-h_0)^2 e^{-2\lambda(Z-h_0)}
= e^{2\lambda h_0}\,(Z-h_0)^2 e^{-2\lambda Z},
\]
so that
\[
\mathbb{E}[Y_\lambda^2]
= 
e^{2\lambda h_0}\mathbb{E}[(Z-h_0)^2 e^{-2\lambda Z}]
=e^{2\lambda h_0}(\mathbb{E}[Z^2 e^{-2\lambda Z}]
-2h_0\,\mathbb{E}[Z e^{-2\lambda Z}]
+h_0^2\,\mathbb{E}[e^{-2\lambda Z}]).
\]
Using the Poisson moment generating function with $t=-2\lambda$ gives
\[
\mathbb{E}[e^{-2\lambda Z}]
=\exp\!\big(\mu(e^{-2\lambda}-1)\big),
\quad
\mathbb{E}[Z e^{-2\lambda Z}]
=\mu e^{-2\lambda}\exp\!\big(\mu(e^{-2\lambda}-1)\big),
\]
and
\[
\mathbb{E}[Z^2 e^{-2\lambda Z}]
=\big(\mu e^{-2\lambda}+\mu^2 e^{-4\lambda}\big)\exp\!\big(\mu(e^{-2\lambda}-1)\big).
\]
Therefore, with $\mu=\mu(\lambda)$,
\begin{align*}
\mathbb{E}[Y_\lambda^2]
&= e^{2\lambda h_0}\exp\!\big(\mu(\lambda)(e^{-2\lambda}-1)\big)
\Big[
\mu(\lambda)e^{-2\lambda}
+\mu(\lambda)^2 e^{-4\lambda}
-2h_0\mu(\lambda)e^{-2\lambda}
+h_0^2
\Big].
\end{align*}
It follows that $\lambda\mapsto \mathbb{E}[Y_\lambda^2]$ is continuous.
Since $\lambda\mapsto \mathbb{E}[Y_\lambda]$ and $\lambda\mapsto \mathbb{E}[Y_\lambda^2]$
are continuous
\[
\lambda\mapsto \mathrm{Var}(Y_\lambda)
=\mathbb{E}[Y_\lambda^2]-\big(\mathbb{E}[Y_\lambda]\big)^2
\]
is continuous.\\

\textbf{Convergence of the estimators of $f(\lambda)$ and $g(\lambda)$}\\



Note that conditioning on $\hat\lambda$, the quantity $\hat\mu=\mu(\hat\lambda)$ is fixed, and
$Z_{\hat\mu}\mid \hat\lambda \sim \mathrm{Poisson}(\hat\mu)$ by definition. Therefore,
\[
\mathbb{E}\!\left[(Z_{\hat{\mu}}-h_0)^2e^{-\hat{\lambda}(Z_{\hat{\mu}}-h_0)}\mid \hat{\lambda}\right]
= f(\hat\lambda),
\]
and similarly
\[
\mathbf{Var}\!\left[(Z_{\hat{\mu}}-h_0)e^{-\hat{\lambda}(Z_{\hat{\mu}}-h_0)}\mid \hat{\lambda}\right]
= g(\hat\lambda).
\]
Since $\hat\lambda\xrightarrow{P}\lambda^{\circ}$ and $f$ is continuous, the continuous mapping theorem yields
\[
f(\hat\lambda)\xrightarrow{P} f(\lambda^{\circ})
=\mathbb{E}\Big[(Z_{\mu}-h_0)^2e^{-\lambda^{\circ}(Z_{\mu}-h_0)}\Big],
\]
where $\mu=\mu(\lambda^{\circ})$. Likewise, continuity of $g$ implies
\[
g(\hat\lambda)\xrightarrow{P} g(\lambda^{\circ})
=\mathbf{Var}\Big((Z_{\mu}-h_0)e^{-\lambda^{\circ}(Z_{\mu}-h_0)}\Big).
\]

\end{proof}

\end{document}
